\chapter{The Critical Line}
\begin{center}\itshape
Riemann guessed that all the zeros lie on one line.\\
Nobody has proved it, but we can prove that infinitely many do,\\
and, if he was right, that they repel each other like the eigenvalues of a random matrix.
\end{center}

Godfrey Harold Hardy visited Harald Bohr in Copenhagen almost every summer, and, as George Pólya remembered it, their programme for the holidays always began with the same first item: prove the Riemann hypothesis. One year Hardy had to go home across the North Sea in a small boat, in bad weather. Before leaving he sent Bohr a postcard: ``I have proved the Riemann hypothesis. G.~H.~Hardy.'' If the boat sank and he drowned, everybody would have to believe that he had proved the Riemann hypothesis, and Hardy, who treated God as a personal opponent, was sure that God would not let him have that honour. The boat did not sink.

What Hardy really proved, in a note of three pages in the \emph{Comptes rendus} of 1914, was the first theorem about zeros \emph{on} the critical line: there are infinitely many of them. Before him, it was not even known that the line carries more than finitely many zeros. His proof never finds a zero. It shows that Riemann's function $\Xi(t)=\xi(\frac12+it)$, which is real for real $t$, cannot keep a constant sign, because its moments against $\cosh$ behave in a way that no eventually positive function can imitate. The secret is again the theta function, the same one that gave the functional equation.

This last chapter is about the critical line. We prove Hardy's theorem, and then the mean value $\int_0^T|\zeta(\frac12+it)|^2\dd t\sim T\log T$ of Hardy and Littlewood (1918), with the help of an elegant inequality of Hilbert. Then we assume the Riemann hypothesis and prove the theorem of Hugh Montgomery (1973) on the pair correlation of the zeros, the theorem whose meeting with random matrix theory over afternoon tea Volume~12 told: we deduce from it that at least two thirds of the zeros are simple.

\begin{reminder}[from Volumes~5,~6 and~12, and Chapter~1]
\emph{The functional equation} (Volume~5, Theorem~12.8). With $\omega(u)=\sum_{n\ge1}e^{-\pi n^2u}$ and $\Lambda_\zeta(s)=\pi^{-s/2}\Gamma(\frac s2)\zeta(s)$,
\[
\Lambda_\zeta(s)=-\frac1s-\frac1{1-s}+\int_1^\infty\omega(u)\bigl(u^{s/2}+u^{(1-s)/2}\bigr)\frac{\mathrm du}u.
\]
\emph{Theta} (Volume~5, Theorem~10.9 and Corollary~10.10). Poisson summation, and $\theta(1/u)=\sqrt u\,\theta(u)$ for $\theta(u)=\sum_{n\in\Z}e^{-\pi n^2u}=1+2\omega(u)$.
\emph{Riemann's $\xi$} (Volume~6, Proposition~6.2). $\xi(s)=\frac12s(s-1)\Lambda_\zeta(s)$ is entire, $\xi(1-s)=\xi(s)$, $\xi(\bar s)=\overline{\xi(s)}$, and its zeros are the non-trivial zeros of $\zeta$, with multiplicity.
\emph{Counting zeros} (Volume~6, Theorem~6.9 and Corollary~6.10). $N(T)=\frac T{2\pi}\log\frac T{2\pi e}+O(\log T)$, and at most $C\log(|u|+2)$ ordinates lie in $[u,u+1]$.
\emph{The explicit formula toolkit} (Volume~6, Lemmas~7.3 and~7.4). Good heights $T'\in[T,T+1]$ with $|\frac{\zeta'}\zeta(\sigma\pm iT')|\le C\log^2T$ for $-1\le\sigma\le3$, and $\frac{\zeta'}\zeta(s)=O(\log(|s|+2))$ for $\operatorname{Re}s\le-1$ away from the trivial zeros.
\emph{Euler--Maclaurin} (Chapter~1, proof of Lemma~\ref{lem:em}). For $\sigma>0$ and an integer $K\ge1$, with $P_1(u)=\{u\}-\frac12$,
\[
\zeta(s)=\sum_{n\le K}n^{-s}+\frac{K^{1-s}}{s-1}-\frac12K^{-s}-s\int_K^\infty P_1(u)u^{-s-1}\dd u.
\]
\end{reminder}

Since $\Lambda$ is von Mangoldt's function in this book, we write $\Lambda_\zeta$ for the completed zeta function, and we work with $\xi$ and
\[
\Xi(t)=\xi\Bigl(\frac12+it\Bigr)\qquad(t\in\C).
\]
By $\xi(1-s)=\xi(s)$ and $\xi(\bar s)=\overline{\xi(s)}$, $\Xi$ is even and real on the real line, and its real zeros are exactly the ordinates of the zeros of $\zeta$ on the critical line.

\section{Hardy's theorem}

The first step writes $\Xi$ as a cosine transform. Put $s=\frac12+it$ in the functional equation. Then $s(1-s)=t^2+\frac14$, and $u^{s/2}+u^{(1-s)/2}=2u^{1/4}\cos(\frac t2\log u)$, so
\[
\begin{aligned}\Xi(t)&=-\frac12\Bigl(t^2+\frac14\Bigr)\Lambda_\zeta(s)\\&=\frac12-\Bigl(t^2+\frac14\Bigr)\\&\quad\times\int_1^\infty\omega(u)u^{-3/4}\cos\Bigl(\frac t2\log u\Bigr)\dd u.\end{aligned}
\]
Substituting $u=e^{2y}$, and using
\[
\int_0^\infty e^{-y/2}\cos(ty)\dd y=\frac{1/2}{1/4+t^2},
\]
gives the first lemma.

\begin{lemma}\label{lem:xicos}
Let $h(y)=e^{-y/2}-2e^{y/2}\omega(e^{2y})$ for $y\ge0$. For real $t$,
\[
\frac{\Xi(t)}{t^2+\frac14}=\int_0^\infty h(y)\cos(ty)\dd y,
\]
and there is a constant $C$ with $|\Xi(t)|\le C(1+|t|)^3e^{-\pi|t|/4}$.
\end{lemma}
\begin{proof}
The formula is the computation above. For the bound, $\pi^{-s/2}\Gamma(\frac s2)$ has modulus $\pi^{-1/4}|\Gamma(\frac14+\frac{it}2)|$, which is at most $C|t|^{-1/4}e^{-\pi|t|/4}$ for $|t|\ge2$ by Stirling's formula on vertical lines (Volume~6, equation~(6.1)), and $|\zeta(\frac12+it)|\le1+3|s|$ for $|t|\ge1$ (Volume~6, proof of Proposition~6.2(c)). Multiply by $\frac12|s(s-1)|\le(1+|t|)^2$; on $|t|\le2$, $\Xi$ is bounded.
\end{proof}

The function $h$ is continuous, and integrable on $[0,\infty[$ since $\omega(e^{2y})\le2e^{-\pi e^{2y}}$. Its even extension $y\mapsto h(|y|)$ is continuous and integrable, and by the lemma its Fourier transform $\xi'\mapsto2\int_0^\infty h(y)\cos(2\pi\xi'y)\dd y$ is $\frac{2\Xi(2\pi\xi')}{(2\pi\xi')^2+1/4}$, which is integrable. Fourier inversion (Volume~5, Theorem~10.6), after the change of variables $t=2\pi\xi'$, gives Hardy's identity: for real $x\ge0$,
\begin{equation}\label{eq:hardyid}
\int_0^\infty\frac{\Xi(t)}{t^2+\frac14}\cos(xt)\dd t=\frac\pi2\Bigl(e^{-x/2}-2e^{x/2}\omega(e^{2x})\Bigr).
\end{equation}
A direct computation verifies it to twenty digits at $x=0$, $\frac12$ and $1$; at $x=0$ both sides equal $1.4350248428\ldots$

Both sides of \eqref{eq:hardyid} make sense for complex $x$. On the left, $|\cos(xt)|\le e^{|\operatorname{Im}x|t}$, and Lemma~\ref{lem:xicos} makes the integral converge, locally uniformly, as long as $|\operatorname{Im}x|<\frac\pi4$. On the right, $\omega(z)$ is holomorphic for $\operatorname{Re}z>0$, and $\operatorname{Re}e^{2x}>0$ exactly when $\cos(2\operatorname{Im}x)>0$. So both sides are holomorphic in the strip $|\operatorname{Im}x|<\frac\pi4$ (Volume~5, Theorem~11.15), agree on $[0,\infty[$, and therefore agree on the whole strip. On the boundary of the strip, at $x=\frac{i\pi}4$, the point $e^{2x}$ reaches $i$, on the boundary of the half-plane where $\theta$ lives. Everything hinges on how $\theta$ behaves there.

\begin{lemma}[theta near $i$]\label{lem:thetai}
For $\operatorname{Re}\delta>0$,
\[
\theta(i+\delta)=\delta^{-1/2}\sum_{m\in\Z}e^{-\pi(m-\frac12)^2/\delta},
\]
with the principal square root. Consequently every derivative of $\theta(i+\delta)$ tends to $0$ when $\delta\to0$ inside a sector $|\arg\delta|\le\frac\pi{12}$.
\end{lemma}
\begin{proof}
Since $e^{-\pi in^2}=(-1)^{n^2}=(-1)^n$, $\theta(i+\delta)=\sum_n(-1)^ne^{-\pi n^2\delta}$. For real $\delta>0$ apply Poisson summation (Volume~5, Theorem~10.9) to $f(u)=e^{-\pi\delta u^2}e^{i\pi u}$, for which $f(n)=(-1)^ne^{-\pi\delta n^2}$. By Volume~5, Theorem~10.3,
$\hat f(\eta)=\int e^{-\pi\delta u^2}e^{-2\pi iu(\eta-\frac12)}\dd u=\delta^{-1/2}e^{-\pi(\eta-\frac12)^2/\delta}$, and both $f$ and $\hat f$ decay fast. This gives the formula for $\delta>0$. Both sides are holomorphic on $\operatorname{Re}\delta>0$, where $\operatorname{Re}\frac1\delta>0$, so they agree there.

Let $K(w)=\theta(i+w)$. Take $\delta$ with $|\arg\delta|\le\frac\pi{12}$, $r=|\delta|\le1$, and a point $w$ with $|w-\delta|\le\frac r2$. Then $w=\delta(1+u)$ with $|u|\le\frac12$, so $|\arg(1+u)|\le\frac\pi6$, $\operatorname{Re}w>0$, $|w|\ge\frac r2$, and $|\arg\frac1w|\le\frac\pi{12}+\frac\pi6=\frac\pi4$. Hence $\operatorname{Re}\frac1w\ge\frac{\cos(\pi/4)}{3r/2}\ge\frac{0.47}r$, and using $(j+\frac12)^2\ge\frac14+j$,
\[
|K(w)|\le\Bigl(\frac2r\Bigr)^{1/2}\cdot2\sum_{j\ge0}e^{-0.47\pi(j+\frac12)^2/r}\le4r^{-1/2}e^{-0.36/r}.
\]
By the Cauchy estimates on the disc of radius $\frac r2$ around $\delta$, $|K^{(j)}(\delta)|\le j!\,(\frac2r)^j\cdot4r^{-1/2}e^{-0.36/r}$, which tends to $0$ as $r\to0$.
\end{proof}

\begin{think}
On the real axis $\theta(u)$ blows up like $u^{-1/2}$ at $0$, because about $u^{-1/2}$ terms are close to $1$. Near $i$ the same terms are $\pm1$ with alternating signs. Before reading the lemma, can you see why the alternating sum is not just bounded but smaller than every power of $|\delta|$?
\end{think}

\begin{theorem}[Hardy, 1914]\label{thm:hardy}
Infinitely many zeros of $\zeta$ lie on the critical line.
\end{theorem}
\begin{proof}
\emph{The moments.} For $0\le\alpha<\frac\pi4$ put $x=i\alpha$ in \eqref{eq:hardyid}. Since $\cos(i\alpha t)=\cosh(\alpha t)$ and $2\omega=\theta-1$,
\[
\Phi(\alpha):=\int_0^\infty\frac{\Xi(t)}{t^2+\frac14}\cosh(\alpha t)\dd t=\frac\pi2\Bigl(2\cos\frac\alpha2-e^{i\alpha/2}\theta(e^{2i\alpha})\Bigr).
\]
Differentiating under the integral (allowed by Lemma~\ref{lem:xicos} for $\alpha\le\alpha_0<\frac\pi4$), for every $n\ge0$
\[
\Phi^{(2n)}(\alpha)=\int_0^\infty\frac{\Xi(t)}{t^2+\frac14}t^{2n}\cosh(\alpha t)\dd t.
\]

\emph{The limit at $\frac\pi4$.} Let $\varepsilon=\frac\pi4-\alpha$. Then $e^{2i\alpha}=ie^{-2i\varepsilon}$, and $\delta(\alpha)=e^{2i\alpha}-i=2\sin\varepsilon\,e^{-i\varepsilon}$: it tends to $0$ inside the sector of Lemma~\ref{lem:thetai}. By the chain rule and induction, the $k$-th derivative of $\alpha\mapsto K(\delta(\alpha))$ is a combination $\sum_{j=1}^kc_{k,j}e^{2ij\alpha}K^{(j)}(\delta(\alpha))$ with constant coefficients, so by the lemma all derivatives of $\theta(e^{2i\alpha})$, and by Leibniz's rule those of $e^{i\alpha/2}\theta(e^{2i\alpha})$, tend to $0$ as $\alpha\to\frac\pi4$. Therefore
\[
\lim_{\alpha\to\pi/4}\Phi^{(2n)}(\alpha)=\frac\pi2\cdot2\Bigl(-\frac14\Bigr)^n\cos\frac\pi8=\frac{(-1)^n\pi\cos\frac\pi8}{4^n}.
\]
In particular, these limits are bounded by $\pi$.

\emph{The contradiction.} Suppose that $\Xi$ has only finitely many real zeros. Then there is $T_0\ge1$ such that $\Xi$ has no zero on $[T_0,\infty[$, so it has a constant sign $\varepsilon_0$ there, and $f(t)=\varepsilon_0\frac{\Xi(t)}{t^2+1/4}$ is positive for $t>T_0$. Split $\varepsilon_0\Phi^{(2n)}(\alpha)$ at $T_0$. The integral over $[0,T_0]$ tends, as $\alpha\to\frac\pi4$, to $\int_0^{T_0}f(t)t^{2n}\cosh\frac{\pi t}4\dd t$, of modulus at most $AT_0^{2n}$ with $A=\int_0^{T_0}|f|\cosh\frac{\pi t}4\dd t$. Hence the integral over $[T_0,\infty[$ has a limit $\ell_n$ with $|\ell_n|\le\pi+AT_0^{2n}$. Its integrand is positive and increases with $\alpha$, so by monotone convergence
\[
\ell_n=\int_{T_0}^\infty f(t)t^{2n}\cosh\frac{\pi t}4\dd t\ge\int_{2T_0}^{2T_0+1}f(t)t^{2n}\cosh\frac{\pi t}4\dd t\ge(2T_0)^{2n}B,
\]
with $B=\int_{2T_0}^{2T_0+1}f(t)\cosh\frac{\pi t}4\dd t>0$. So $4^nB\le\pi T_0^{-2n}+A\le\pi+A$ for every $n$, which is absurd.
\end{proof}

\begin{beauty}
The proof never sees a zero. It weighs $\Xi$ against $t^{2n}\cosh\frac{\pi t}4$, the heaviest weight it can bear, and finds that the answers stay bounded while the weights explode. An eventually positive function cannot do that: its far tail would dominate. Only a function that keeps changing sign, so that its tail cancels itself, can have such tame moments. The cancellation comes from $\theta$ at the point $i$, where the transformation $\theta(1/u)=\sqrt u\,\theta(u)$ turns an alternating sum into a series smaller than every power: the same symmetry that gives the functional equation forces zeros onto the line.
\end{beauty}

Hardy's theorem says nothing about the proportion of zeros on the line. Hardy and Littlewood showed in 1921 that $[0,T]$ contains at least $cT$ of their ordinates, still a vanishing proportion of the $N(T)\approx\frac T{2\pi}\log T$ zeros. In 1942 Atle Selberg, working almost alone in occupied Norway, found the decisive idea: mollify $\zeta$ before measuring its sign changes, as in the zero density estimates of Chapter~1.

\begin{theorem}[Selberg, 1942; Levinson, 1974; Conrey, 1989; Pratt, Robles, Zaharescu and Zeindler, 2020; quoted]\label{thm:proportion}
A positive proportion of the zeros of $\zeta$ lie on the critical line: Selberg proved $\liminf N_0(T)/N(T)>0$, where $N_0(T)$ counts the zeros $\frac12+i\gamma$ with $0<\gamma\le T$. Norman Levinson proved that the proportion exceeds $\frac13$, Brian Conrey that it exceeds $\frac25$, and Pratt, Robles, Zaharescu and Zeindler that it exceeds $0.4172$, more than five twelfths.
\end{theorem}

On the computational side, David Platt and Timothy Trudgian verified in 2021 that every zero with $0<\gamma\le3\cdot10^{12}$ lies on the line and is simple, by the method of Turing recalled in Volume~6, Chapter~6.

\section{The mean square on the critical line}

How big is $\zeta$ on the critical line, on average? For $\sigma>1$, $\zeta(\sigma+it)$ is a convergent Dirichlet series, and its mean square is $\sum n^{-2\sigma}=\zeta(2\sigma)$ (Exercise~\ref{ex:meansqsigma} extends this to $\sigma>\frac12$). At $\sigma=\frac12$ the diagonal $\sum n^{-1}$ diverges, and the answer should be a logarithm. To get it we need $\zeta(\frac12+it)$ as a Dirichlet polynomial of length about $t$, not $t^{1+\varepsilon}$ as in Chapter~1, and a mean value theorem without the logarithmic loss of Lemma~\ref{lem:mvt}.

\begin{lemma}[first derivative test for integrals]\label{lem:firstderiv}
Let $F$ be real and twice continuously differentiable on $[a,b]$, with $F'$ monotone, of constant sign, and $|F'|\ge m>0$. Let $g$ be positive, decreasing and continuously differentiable. Then
\[
\Bigl|\int_a^bg(u)e(F(u))\dd u\Bigr|\le\frac{g(a)}m.
\]
\end{lemma}
\begin{proof}
Write $e(F)=\frac{(e(F))'}{2\pi iF'}$ and integrate by parts:
\[
\int_a^bg\,e(F)=\Bigl[\frac{g\,e(F)}{2\pi iF'}\Bigr]_a^b-\frac1{2\pi i}\int_a^be(F)\Bigl(\frac g{F'}\Bigr)'.
\]
The bracket is at most $\frac{g(a)+g(b)}{2\pi m}\le\frac{g(a)}{\pi m}$. Next $(\frac g{F'})'=\frac{g'}{F'}-\frac{gF''}{F'^2}$, and $\int_a^b\frac{|g'|}{|F'|}\le\frac{g(a)-g(b)}m\le\frac{g(a)}m$. Since $F'$ is monotone, $F''$ has a constant sign, and $\int_a^b\frac{g|F''|}{F'^2}\le g(a)\bigl|\frac1{F'(a)}-\frac1{F'(b)}\bigr|\le\frac{g(a)}m$, the two numbers having the same sign and modulus at most $\frac1m$. Altogether the integral is at most $\frac{2g(a)}{\pi m}<\frac{g(a)}m$.
\end{proof}

\begin{lemma}[$\zeta$ as a sum of length $t$]\label{lem:zetalengtht}
Let $0<\sigma\le1$, $t\ge0$, $s=\sigma+it$ and $x\ge t+3$. Then
\[
\zeta(s)=\sum_{n\le x}n^{-s}+\frac{x^{1-s}}{s-1}+E,\qquad|E|\le4x^{-\sigma}.
\]
\end{lemma}
\begin{proof}
Let $K=\lfloor x\rfloor\ge t+2$. Let $P_2$ be the periodic function of Chapter~1, $P_2(u)=\frac{\{u\}^2-\{u\}}2+\frac1{12}$, with $P_2'=P_1$ away from the integers and $-\frac1{24}\le P_2\le\frac1{12}$. Two integrations by parts give its Fourier coefficients: $\int_0^1P_2(u)e(-ku)\dd u=\frac1{4\pi^2k^2}$ for $k\ne0$, and $0$ for $k=0$. They are summable, so (Volume~5, Corollary~8.10)
\[
P_2(u)=\sum_{k\ge1}\frac{\cos2\pi ku}{2\pi^2k^2}.
\]
Integrating by parts in the reminder formula,
\[
\begin{aligned}\int_K^\infty P_1u^{-s-1}\dd u&=-\frac{K^{-s-1}}{12}\eqbr+(s+1)\int_K^\infty P_2u^{-s-2}\dd u,\end{aligned}
\]
and the series of $P_2$ can be integrated term by term since $\sum k^{-2}\int_K^\infty u^{-\sigma-2}\dd u<\infty$:
\[
\begin{aligned}s\int_K^\infty P_1(u)u^{-s-1}\dd u&=-\frac{sK^{-s-1}}{12}\eqbr+s(s+1)\sum_{k\ge1}\frac1{2\pi^2k^2}\int_K^\infty\cos(2\pi ku)\,u^{-s-2}\dd u.\end{aligned}
\]
Write $\cos(2\pi ku)u^{-it}=\frac12\bigl(e(F_+(u))+e(F_-(u))\bigr)$ with $F_\pm(u)=\pm ku-\frac t{2\pi}\log u$. Then $F_\pm'(u)=\pm k-\frac t{2\pi u}$ is increasing, and for $u\ge K\ge t$, $|F_+'|\ge k-\frac1{2\pi}\ge\frac k2$ and $|F_-'|\ge k$. Lemma~\ref{lem:firstderiv} with $g(u)=u^{-\sigma-2}$, on $[K,b]$ with $b\to\infty$, bounds the $k$-th integral by $\frac12(\frac2k+\frac1k)K^{-\sigma-2}$. Since $|s|\le t+1\le K$ and $|s+1|\le t+2\le K$, the sum over $k$ is at most $K^2\cdot\frac{3\zeta(3)}{4\pi^2}K^{-\sigma-2}\le0.1K^{-\sigma}$, and the first term at most $\frac1{12}K^{-\sigma}$.

Finally $\frac{K^{1-s}-x^{1-s}}{s-1}=-\int_K^xu^{-s}\dd u$ has modulus at most $K^{-\sigma}$, and $\sum_{n\le x}=\sum_{n\le K}$. Collecting, $|E|\le(1+\frac12+\frac1{12}+0.1)K^{-\sigma}\le1.7\cdot2^\sigma x^{-\sigma}$, as $K\ge x-1\ge\frac x2$.
\end{proof}

The tail $\sum_{x<n}n^{-s}$ has been replaced by the integral $\int_x^\infty u^{-s}\dd u$: as long as the phase $t\log n$ moves by less than $2\pi$ from one integer to the next, the sum cannot tell itself apart from the integral. Now the mean value. The off-diagonal terms of $|\sum a_nn^{-it}|^2$ involve $\frac1{\log(n/m)}$, and the key is an inequality about the matrix $\frac1{m-n}$.

\begin{theorem}[Hilbert's inequality, with Schur's constant]\label{thm:hilbert}
For all complex numbers $x_1,\dots,x_N$,
\[
\Bigl|\sum_{\substack{m,n\le N\\m\ne n}}\frac{x_m\overline{x_n}}{m-n}\Bigr|\le\pi\sum_{n\le N}|x_n|^2.
\]
\end{theorem}
\begin{proof}
Let $P(u)=\sum_nx_ne(nu)$. For an integer $k\ne0$, $\int_0^1ue(ku)\dd u=\frac1{2\pi ik}$ (integrate by parts) and $\int_0^1e(ku)\dd u=0$, while $\int_0^1(u-\frac12)\dd u=0$. Expanding $|P|^2$,
\[
\int_0^1\Bigl(u-\frac12\Bigr)|P(u)|^2\dd u=\frac1{2\pi i}\sum_{m\ne n}\frac{x_m\overline{x_n}}{m-n}.
\]
The left side has modulus at most $\frac12\int_0^1|P|^2=\frac12\sum|x_n|^2$.
\end{proof}

The inequality goes back to David Hilbert's lectures on integral equations and was first published, with the constant $2\pi$, in Hermann Weyl's thesis of 1908; Issai Schur found the constant $\pi$ in 1911, and it cannot be improved (Exercise~\ref{ex:hilbertopt}). Hugh Montgomery and Robert Vaughan extended it in 1974 to arbitrary well-spaced real numbers in place of the integers, which gives the sharp mean value theorem alluded to in Chapter~1. For our purpose a trick suffices.

\begin{proposition}\label{prop:meansqpoly}
Let $N\ge1$, $T>0$ and $S(t)=\sum_{n\le N}n^{-1/2-it}$. Then
\[
\Bigl|\int_0^T|S(t)|^2\dd t-T\sum_{n\le N}\frac1n\Bigr|\le16N.
\]
\end{proposition}
\begin{proof}
Expanding, $|S(t)|^2=\sum_{m,n}(mn)^{-1/2}(\frac nm)^{it}$. The diagonal gives $T\sum\frac1n$, and the rest is
\[
\sum_{m\ne n}\frac{(mn)^{-1/2}}{i\log(n/m)}\Bigl(\Bigl(\frac nm\Bigr)^{iT}-1\Bigr).
\]
The part with $-1$ vanishes: exchanging $m$ and $n$ changes the sign of each term. For the other part, write $u=\log\frac nm$ and
\[
\frac1u=\frac1{2\sinh(u/2)}+k(u),\qquad k(u)=\frac1u-\frac1{2\sinh(u/2)},\qquad2\sinh\frac u2=\frac{n-m}{\sqrt{mn}}.
\]
The first piece contributes $\frac1i\sum_{m\ne n}\frac{n^{iT}\overline{m^{iT}}}{n-m}$, at most $\pi N$ by Theorem~\ref{thm:hilbert} with $x_n=n^{iT}$. For the second, $k$ is odd, $0\le k(u)\le\frac1u$ for $u>0$ (as $2\sinh\frac u2\ge u$), and $k(u)\le\frac u{20}$ for $0<u\le1$ (by $\sinh y-y\le\frac{y^3}6\cosh y$). So the second piece is at most $2\sum_{m<n}(mn)^{-1/2}k(\log\frac nm)$. When $m<n\le2m$, $\log\frac nm\le\frac{n-m}m$ and the term is at most $\frac{n-m}{20m^2}$; summing over $n=m+r$, $r\le m$, gives at most $\frac1{20}$ for each $m$. When $n>2m$, $k\le\frac1{\log2}$, and these terms add up to at most $\frac1{\log2}(\sum_{m\le N}m^{-1/2})^2\le\frac{4N}{\log2}$. The second piece is therefore at most $2(\frac N{20}+5.78N)\le12N$, and $\pi+12<16$.
\end{proof}

\begin{theorem}[Hardy and Littlewood, 1918]\label{thm:hardylittlewood}
As $T\to\infty$,
\[
\int_0^T\Bigl|\zeta\Bigl(\frac12+it\Bigr)\Bigr|^2\dd t=T\log T+O\bigl(T\sqrt{\log T}\bigr).
\]
\end{theorem}
\begin{proof}
Let $x=T+3$ and $N=\lfloor x\rfloor$. For $0\le t\le T$, Lemma~\ref{lem:zetalengtht} gives $\zeta(\frac12+it)=S(t)+r(t)$, with $S$ as in Proposition~\ref{prop:meansqpoly} and
\[
|r(t)|\le\frac{\sqrt x}{|{-\frac12}+it|}+\frac4{\sqrt x}\le\frac{\sqrt x}{\max(\frac12,t)}+\frac4{\sqrt x}.
\]
So $\int_0^T|r|^2\le2x(4+\int_1^\infty t^{-2}\dd t)+\frac{32T}x=O(T)$. By the proposition and $\sum_{n\le N}\frac1n=\log N+O(1)$, $\int_0^T|S|^2=T\log T+O(T)$. By Cauchy--Schwarz,
\[
\Bigl|\int_0^T|S+r|^2-\int_0^T|S|^2\Bigr|\le2\Bigl(\int_0^T|S|^2\Bigr)^{1/2}\Bigl(\int_0^T|r|^2\Bigr)^{1/2}+\int_0^T|r|^2=O\bigl(T\sqrt{\log T}\bigr).
\]
\end{proof}

Albert Ingham found the second term in 1926: the integral is $T\log\frac T{2\pi}+(2\gamma-1)T+O(T^{1/2}\log T)$, with Euler's constant $\gamma$. Figure~\ref{fig:meansq} compares it with the truth. The fourth moment $\frac1{2\pi^2}T\log^4T$ is also Ingham's (1926); no higher moment is known, and the conjectured constants come from random matrices (Volume~8, Chapter~13).

\begin{figure}[ht]
\centering
\fig{ch14-meansq}
\caption{The mean square $\frac1T\int_0^T|\zeta(\frac12+it)|^2\dd t$ minus $\log\frac T{2\pi}$, for $10\le T\le2000$ (charcoal), computed numerically with the Riemann--Siegel formula, against Ingham's constant $2\gamma-1=0.1544$ (rust). Each zero of $\zeta$ is a point where $|\zeta|^2$ vanishes, and the curve wiggles at the scale of their spacing.}
\label{fig:meansq}
\end{figure}

\emph{Two sums of length $\sqrt t$.} Lemma~\ref{lem:zetalengtht} needs $t$ terms. What becomes of the terms with $\sqrt t<n\le t$? Their phases $t\log n$ have a stationary point with respect to Poisson summation exactly when $\frac t{2\pi n}$ is an integer $m$, and the sum over these points rebuilds a second, dual Dirichlet polynomial in $m\le\frac t{2\pi n}$. The result is the approximate functional equation.

\begin{theorem}[Hardy and Littlewood, 1921; quoted]\label{thm:afe}
Let $0<\sigma<1$, and $x,y\ge1$ with $2\pi xy=t$. With $\chi(s)=\pi^{s-1/2}\frac{\Gamma((1-s)/2)}{\Gamma(s/2)}$, so that $\zeta(s)=\chi(s)\zeta(1-s)$,
\[
\zeta(s)=\sum_{n\le x}n^{-s}+\chi(s)\sum_{n\le y}n^{s-1}+O\bigl(x^{-\sigma}+t^{1/2-\sigma}y^{\sigma-1}\bigr).
\]
\end{theorem}

Riemann himself knew much more. In 1932 Carl Ludwig Siegel deciphered Riemann's unpublished notes in the Göttingen library and found an asymptotic expansion that Riemann had used to compute the first zeros by hand.

\begin{theorem}[Riemann, Siegel, 1932; quoted]\label{thm:riemannsiegel}
Let $Z(t)=e^{i\vartheta(t)}\zeta(\frac12+it)$ be Hardy's function, real for real $t$ (Exercise~\ref{ex:zreal}). For $t\ge2\pi$, let $a=\sqrt{t/2\pi}$, $N=\lfloor a\rfloor$ and $p=a-N$. Then
\[
Z(t)=2\sum_{n\le N}\frac{\cos(\vartheta(t)-t\log n)}{\sqrt n}+(-1)^{N-1}a^{-1/2}\Bigl(\frac{\cos2\pi(p^2-p-\frac1{16})}{\cos2\pi p}+O(t^{-1/2})\Bigr).
\]
\end{theorem}

This is how zeros are computed: $\sqrt{t/2\pi}$ cosines instead of the $t$ terms of Lemma~\ref{lem:zetalengtht}. Andrew Odlyzko and Arnold Schönhage showed in 1988 how to evaluate the main sum at many points $t$ at once with the fast Fourier transform, at an average cost of $t^{\varepsilon}$ operations per value, and this is what made the computation of billions of zeros possible. Direct computation gives all $10\,142$ zeros with $0<\gamma\le10^4$ in about a second.

\section{Montgomery's pair correlation}

From now on we assume the Riemann hypothesis, and write the zeros as $\frac12+i\gamma$. Volume~8, Chapter~13 told how Montgomery, studying the differences $\gamma-\gamma'$, found the function $1-(\frac{\sin\pi u}{\pi u})^2$, and how Freeman Dyson recognised it at tea as the pair correlation of the eigenvalues of random matrices. Here is the theorem behind the story. For $T\ge2$ and real $\alpha$, let
\[
\begin{gathered}F(\alpha)=\frac{2\pi}{T\log T}\sum_{0<\gamma,\gamma'\le T}T^{i\alpha(\gamma-\gamma')}w(\gamma-\gamma'),\\w(u)=\frac4{4+u^2},\end{gathered}
\]
the sum running over pairs of zeros counted with multiplicity. The scale $T^{i\alpha(\gamma-\gamma')}=e^{i\alpha(\gamma-\gamma')\log T}$ measures differences in units of $\frac{2\pi}{\log T}$, the mean spacing of the zeros up to height $T$, and $\alpha$ is the dual frequency. The weight $w$ only makes the sum converge.

The tool is an explicit formula in which the zeros come with the weight $\frac1{1+(t-\gamma)^2}$. Its sum over zeros mixes two zeros at once when squared, and this weight was chosen so that it squares well.

\begin{lemma}\label{lem:cauchykernels}
For real $a,b$,
\[
\int_{-\infty}^\infty\frac{\mathrm dt}{(1+(t-a)^2)(1+(t-b)^2)}=\frac{2\pi}{4+(a-b)^2}.
\]
\end{lemma}
\begin{proof}
Close the contour in the upper half-plane. The poles there are $a+i$ and $b+i$, with residues $\frac1{2i(1+(d+i)^2)}$ and $\frac1{2i(1+(-d+i)^2)}$, where $d=a-b$. Since $1+(d+i)^2=d(d+2i)$ and $1+(-d+i)^2=d(d-2i)$, their sum is $\frac1{2id}\bigl(\frac1{d+2i}+\frac1{d-2i}\bigr)=\frac1{i(d^2+4)}$ (for $d=0$ by continuity), and $2\pi i$ times it is the claim.
\end{proof}

\begin{lemma}[Montgomery's explicit formula]\label{lem:montexplicit}
Assume the Riemann hypothesis. For $x\ge1$ and real $t$, with $\tau=|t|+2$,
\[
\begin{aligned}2\sum_\gamma\frac{x^{i\gamma}}{1+(t-\gamma)^2}&=-x^{-1/2}\sum_{n\le x}\Lambda(n)\Bigl(\frac xn\Bigr)^{-1/2+it}\\&\quad-x^{-1/2}\sum_{n>x}\Lambda(n)\Bigl(\frac xn\Bigr)^{3/2+it}\\&\quad+x^{-1+it}\log\tau+O\bigl(x^{1/2}\tau^{-2}+x^{-1}\bigr),\end{aligned}
\]
the sum running over all zeros, with multiplicity. The implied constant is absolute.
\end{lemma}
\begin{proof}
Let $s=\frac12+it$ and $g(w)=\frac{x^{w-1/2}}{(1+w-s)(1+s-w)}$, with simple poles at $s-1$ and $s+1$. Let $c=\frac54$ and $\mathcal F(w)=-\frac{\zeta'}\zeta(w)g(w)$.

\emph{The prime side.} On $\operatorname{Re}w=c$, $-\frac{\zeta'}\zeta(w)=\sum\Lambda(n)n^{-w}$ converges absolutely, and $|g(w)|=O(|w|^{-2})$, so we may integrate term by term. With $u=w-s$, on the line $\operatorname{Re}u=\frac34$,
\[
\frac1{2\pi i}\int_{(c)}n^{-w}g(w)\dd w=x^{-1/2}\Bigl(\frac xn\Bigr)^s\cdot\frac1{2\pi i}\int_{(3/4)}\frac{(x/n)^u}{(1+u)(1-u)}\dd u.
\]
For $y=\frac xn\ge1$, closing to the left (where $|y^u|\le y^{3/4}$ and the integrand is $O(|u|^{-2})$) picks up the pole $u=-1$, with residue $\frac{y^{-1}}2$; for $y\le1$, closing to the right picks up $u=1$, with residue $-\frac y2$, counted with a minus sign. So the last integral is $\frac12\min(y,y^{-1})$, and
\[
\frac1{2\pi i}\int_{(c)}\mathcal F(w)\dd w=\frac{x^{-1/2}}2\Bigl(\sum_{n\le x}\Lambda(n)\Bigl(\frac xn\Bigr)^{s-1}+\sum_{n>x}\Lambda(n)\Bigl(\frac xn\Bigr)^{s+1}\Bigr).
\]

\emph{The residues.} Move the line to the left as in the proof of the explicit formula (Volume~6, Theorem~7.5), over the rectangle with vertices $c\pm iT'$, $-U\pm iT'$, where $T'\ge2\tau$ is a good height and $U=2K+1$. On the horizontal sides and on the left side $|w|\ge2\tau$, so $|1+w-s|\ge\frac{|w|}2$ and $|1+s-w|\ge\frac{|w|}2$, and $|g(w)|\le\frac{4x^{\operatorname{Re}w-1/2}}{|w|^2}$. With the bounds for $\frac{\zeta'}\zeta$ recalled above, the left side contributes $O(T'\log(U+T')U^{-2})$ (here $x\ge1$), which tends to $0$ as $U\to\infty$, and the horizontal sides contribute $O(x^{3/4}\frac{\log^2T'}{T'})$, which tends to $0$ as $T'\to\infty$. The poles inside are: $w=1$, where $-\frac{\zeta'}\zeta$ has residue $1$, giving $g(1)$; each zero $\rho$, giving $-g(\rho)$ for each unit of multiplicity; $w=s-1=-\frac12+it$, where the residue of $g$ is $\frac{x^{s-3/2}}2$, giving $-\frac12x^{-1+it}\frac{\zeta'}\zeta(-\frac12+it)$; and the trivial zeros $-2k$, giving $-g(-2k)$. The pole $s+1$ lies to the right of the line. Under the Riemann hypothesis $g(\rho)=\frac{x^{i\gamma}}{1+(\gamma-t)^2}$, and this series converges absolutely. Hence
\[
\frac1{2\pi i}\int_{(c)}\mathcal F=g(1)-\sum_\gamma\frac{x^{i\gamma}}{1+(t-\gamma)^2}-\frac{x^{-1+it}}2\frac{\zeta'}\zeta\Bigl(-\frac12+it\Bigr)-\sum_{k\ge1}g(-2k).
\]

\emph{The terms.} $|g(1)|=\frac{x^{1/2}}{|\frac32-it||\frac12+it|}=O(x^{1/2}\tau^{-2})$, and $|g(-2k)|\le\frac{x^{-2k-1/2}}{(2k-\frac12)(2k+\frac32)}$, so the trivial zeros contribute $O(x^{-5/2})$. Finally, the logarithmic derivative of the functional equation (Volume~6, proof of Lemma~7.4) is
\[
\frac{\zeta'}\zeta(w)=\log\pi-\frac12\psi_\Gamma\Bigl(\frac w2\Bigr)-\frac12\psi_\Gamma\Bigl(\frac{1-w}2\Bigr)-\frac{\zeta'}\zeta(1-w).
\]
At $w=-\frac12+it$: $\frac{\zeta'}\zeta(\frac32-it)$ is bounded; $\psi_\Gamma(\frac34-\frac{it}2)=\log\tau+O(1)$ by Stirling (Volume~6, Theorem~6.1(c)); and $\psi_\Gamma(-\frac14+\frac{it}2)=\psi_\Gamma(\frac54-\frac{it}2)-\pi\cot\pi(-\frac14+\frac{it}2)=\log\tau+O(1)$, the cotangent being bounded at distance $\frac14$ from the integers. So $\frac{\zeta'}\zeta(-\frac12+it)=-\log\tau+O(1)$. Multiply the identity by $2$ and rearrange.
\end{proof}

A direct computation verifies the lemma with the zeros up to height $10^4$ and the primes up to $10^7$: at $x=10$, $t=100$ both sides equal $0.8285+1.3374i$ to four digits.

Write the lemma as $L(t)=U(t)+V(t)+E(t)$, with $L(t)=2\sum_\gamma\frac{x^{i\gamma}}{1+(t-\gamma)^2}$,
\[
U(t)=-x^{it-1/2}D(t),\quad D(t)=\sum_{n\ge1}a_nn^{-it},\quad a_n=\begin{cases}\Lambda(n)(n/x)^{1/2}&(n\le x),\\\Lambda(n)(x/n)^{3/2}&(n>x),\end{cases}
\]
$V(t)=x^{-1+it}\log(|t|+2)$, and $E(t)=O(x^{1/2}\tau^{-2}+x^{-1})$. Two sums of primes appear.

\begin{lemma}\label{lem:lambdasq}
For $x\ge1$, $\sum_na_n^2=x\log x+O(x)$. Moreover, with
\[
b_n=\log n\cdot\min\Bigl(\Bigl(\frac nx\Bigr)^{1/2},\Bigl(\frac xn\Bigr)^{3/2}\Bigr),
\]
$\sum_nb_n=O(x\log2x)$ and $\sum_nn\log(2n)\,b_n^2=O(x^2\log^32x)$.
\end{lemma}
\begin{proof}
Let $S(y)=\sum_{n\le y}\Lambda(n)^2$. The prime powers $p^k$ with $k\ge2$ contribute $O(\sqrt y\log^2y)$ to $S(y)$ and to $\sum_{n\le y}\Lambda(n)\log n$, so by partial summation and the prime number theorem with error term (Volume~6, Theorem~7.8), $S(y)=\psi(y)\log y-\int_2^y\frac{\psi(u)}u\dd u+O(\sqrt y\log^2y)=y\log y+O(y)$. Partial summation again gives
\[
\begin{aligned}\sum_{n\le x}\Lambda(n)^2n&=xS(x)-\int_1^xS(u)\dd u\\&=\frac{x^2\log x}2+O(x^2),\end{aligned}
\]
\[
\begin{aligned}\sum_{n>x}\frac{\Lambda(n)^2}{n^3}&=-\frac{S(x)}{x^3}+3\int_x^\infty\frac{S(u)}{u^4}\dd u\\&=\frac{\log x}{2x^2}+O(x^{-2}),\end{aligned}
\]
and so
\[
\sum_na_n^2=\frac1x\sum_{n\le x}\Lambda(n)^2n+x^3\sum_{n>x}\frac{\Lambda(n)^2}{n^3}=x\log x+O(x).
\] The sums over $b_n$ are elementary: split at $n=x$ and compare with integrals.
\end{proof}

\begin{theorem}[Montgomery, 1973]\label{thm:montgomerypc}
Assume the Riemann hypothesis. Then $F$ is real, even and non-negative, and for every $\varepsilon>0$, uniformly for $0\le\alpha\le1-\varepsilon$, as $T\to\infty$,
\[
F(\alpha)=\bigl(1+o(1)\bigr)T^{-2\alpha}\log T+\alpha+o(1)
\]
\end{theorem}
\begin{proof}
Let $x=T^\alpha$, so $1\le x\le T^{1-\varepsilon}$, and assume $T\ge3$.

\emph{$F$ as a mean square.} Let
\[
L_0(t)=2\sum_{0<\gamma\le T}\frac{x^{i\gamma}}{1+(t-\gamma)^2}.
\]
Expanding $|L_0|^2$ and using Lemma~\ref{lem:cauchykernels},
\[
\begin{aligned}\int_{-\infty}^\infty|L_0(t)|^2\dd t&=4\sum_{0<\gamma,\gamma'\le T}x^{i(\gamma-\gamma')}\frac{2\pi}{4+(\gamma-\gamma')^2}\\&=2\pi\sum_{0<\gamma,\gamma'\le T}x^{i(\gamma-\gamma')}w(\gamma-\gamma')=F(\alpha)\,T\log T.\end{aligned}
\]
This shows that $F(\alpha)$ is real and non-negative; it is even since $x^{-i(\gamma-\gamma')}$ is the conjugate of $x^{i(\gamma-\gamma')}$.

\emph{From $L_0$ to $L$.} By Volume~6, Corollary~6.10, grouping the ordinates in intervals of length $1$, $\sum_\gamma\frac1{1+(t-\gamma)^2}=O(\log(|t|+2))$, $\sum_{\gamma>T}\frac1{1+(t-\gamma)^2}=O(\frac{\log T}{1+T-t})$ and $\sum_{\gamma\le0}\frac1{1+(t-\gamma)^2}=O(\frac{\log T}{1+t})$ for $0\le t\le T$. So on $[0,T]$, $|L_0|=O(\log T)$ and $|L-L_0|\le e(t)$ with $\int_0^Te=O(\log^2T)$ and $\int_0^Te^2=O(\log^2T)$; hence $\int_0^T|L|^2=\int_0^T|L_0|^2+O(\log^3T)$. Outside $[0,T]$, the same counting gives $|L_0(t)|=O(\frac{\log(|t|+2)}{1+|t|})$ for $t<0$ and $O(\frac{\log T}{1+t-T})$ for $t>T$, so $\int_{\R\setminus[0,T]}|L_0|^2=O(\log^2T)$. Altogether
\[
F(\alpha)\,T\log T=\int_0^T|L(t)|^2\dd t+O(\log^3T).
\]

\emph{The diagonal.} Since $\sum|a_n|<\infty$, $|D|^2$ can be expanded and integrated term by term, and $|\int_0^T(n/m)^{it}\dd t|\le\frac2{|\log(n/m)|}$ for $m\ne n$. So
\[
\int_0^T|D|^2=T\sum a_n^2+O\Bigl(\sum_{m\ne n}\frac{a_ma_n}{|\log(n/m)|}\Bigr).
\]
Bound $a_ma_n\le\frac12(b_m^2+b_n^2)$ when $\frac12<\frac nm<2$, where $|\log\frac nm|\ge\frac{|n-m|}{2m}$ and the sum over $n$ of $\frac1{|\log(n/m)|}$ is at most $4m(1+\log2m)$; and $a_ma_n\le b_mb_n$, $|\log\frac nm|\ge\log2$ otherwise. By Lemma~\ref{lem:lambdasq} the off-diagonal is $O(x^2\log^32x+x^2\log^22x)$, and
\[
\int_0^T|U|^2=x^{-1}\int_0^T|D|^2=T\log x+O(T+x\log^32x).
\]

\emph{The other terms.} $\int_0^T|V|^2=x^{-2}\int_0^T\log^2(t+2)\dd t=x^{-2}(T\log^2T+O(T\log T))$. For the cross term, $\int_0^TU\overline V=-x^{-3/2}\sum_na_n\int_0^Tn^{-it}\log(t+2)\dd t$; here $a_1=0$, and for $n\ge2$ an integration by parts bounds the inner integral by $\frac{3\log(T+2)}{\log2}$. By Lemma~\ref{lem:lambdasq} the cross term is $O(x^{-1/2}\log2x\log T)=O(\log T)$. Finally $\int_0^T|E|^2=O(x+Tx^{-2})$.

\emph{Together.} So $\int_0^T|U+V|^2=T\log x+x^{-2}T\log^2T+O(T+x\log^32x+x^{-2}T\log T)$, and by Cauchy--Schwarz the term $E$ changes this by
\[
O\Bigl(x+Tx^{-2}+\sqrt{(x+Tx^{-2})(T\log T+x^{-2}T\log^2T)}\Bigr),
\]
which is $O\bigl(T^{1-\varepsilon/2}\log T+Tx^{-1}\sqrt{\log T}+Tx^{-2}\log T\bigr)$.
Since $x\le T^{1-\varepsilon}$, every error is $o(T\log T)$ except $O(Tx^{-2}\log T)$, which is $O(\frac1{\log T})$ times $x^{-2}T\log^2T$. Dividing by $T\log T$, with $\log x=\alpha\log T$ and $x^{-2}=T^{-2\alpha}$, gives the theorem.
\end{proof}

For small $\alpha$ the term $T^{-2\alpha}\log T$ dominates: it is the trace of the obvious correlation, each zero with itself counted on the scale $\frac{2\pi}{\log T}$, and its integral over $\alpha$ is $1$. For larger $\alpha$ what remains is $\alpha$. Montgomery conjectured that $F(\alpha)=1+o(1)$ for $\alpha\ge1$, uniformly on bounded intervals, and Fourier inversion turns ``$|\alpha|$ for $|\alpha|\le1$, then $1$'' into the density $1-(\frac{\sin\pi u}{\pi u})^2$ of the beauty box of Volume~8, Chapter~13. Theorem~\ref{thm:montgomerypc} is the statement quoted there as Theorem~13.13. Figure~\ref{fig:falpha} shows $F$ computed from the zeros, and Figure~\ref{fig:paircorr} the pair correlation itself.

\begin{figure}[ht]
\centering
\fig{ch14-falpha}
\caption{Montgomery's function at height $T=5000$, computed from the $4\,520$ zeros with $0<\gamma\le5000$ unfolded to mean spacing $1$ (charcoal), against $\alpha+\ell e^{-2\alpha\ell}$ with $\ell=\log\frac T{2\pi}$ for $\alpha<1$, and Montgomery's conjecture $F=1$ beyond (rust, dashed). At finite height the unfolded version is the fair comparison: the normalisation $\frac{T\log T}{2\pi}$ of the theorem differs from $N(T)$ by the factor $\log\frac T{2\pi e}/\log T$, which tends to $1$ very slowly. The theorem covers only $\alpha<1$; beyond, the computed values follow the conjecture.}
\label{fig:falpha}
\end{figure}

\begin{figure}[ht]
\centering
\fig{ch14-paircorr}
\caption{Pair correlation of the zeros with $10^3<\gamma\le10^4$: the histogram of the differences $\gamma'-\gamma$ of all pairs with $\gamma<\gamma'$, rescaled by the local mean spacing $\frac{2\pi}{\log(\gamma/2\pi)}$, normalised by the number of zeros, against Montgomery's prediction $1-(\frac{\sin\pi u}{\pi u})^2$ (rust) and the constant $1$ of independent points (dashed).}
\label{fig:paircorr}
\end{figure}

\begin{beauty}
In 1973 nobody could compute what $F$ does beyond $\alpha=1$; the explicit formula sees only primes up to $x=T^\alpha$, and at $\alpha=1$ the primes up to $T$ start to interfere with each other in ways that nobody controls. Daniel Goldston and Montgomery proved in 1987 that, under the Riemann hypothesis, Montgomery's conjecture is \emph{equivalent} to a statement about primes: the variance of the number of primes in short intervals, $\int_1^X(\psi(x+\delta x)-\psi(x)-\delta x)^2\dd x\sim\frac12\delta X^2\log\frac1\delta$ in a suitable range of $\delta$. The repulsion between zeros and the randomness of primes in short intervals are the same phenomenon seen from the two sides of the explicit formula.
\end{beauty}

Even the range $\alpha<1$ says something concrete about the zeros: they cannot pile up.

\begin{corollary}[Montgomery, 1973]\label{cor:simplezeros}
Assume the Riemann hypothesis. Then $\sum_{0<\gamma\le T}m_\rho\le(\frac43+o(1))N(T)$, where $m_\rho$ is the multiplicity of $\rho$ and the sum counts each zero with multiplicity. Consequently, at least $(\frac23-o(1))N(T)$ of the zeros with $0<\gamma\le T$ are simple.
\end{corollary}
\begin{proof}
Let $0<\beta<1$ and $r(\alpha)=(1-\frac{|\alpha|}\beta)_+$. Its transform is non-negative (Exercise~\ref{ex:fejervalue}):
\[
\int_{-\beta}^\beta r(\alpha)e^{i\alpha v}\dd\alpha=\beta\Bigl(\frac{\sin(\beta v/2)}{\beta v/2}\Bigr)^2\ge0.
\]
Multiply the definition of $F$ by $r(\alpha)$ and integrate:
\[
\frac{T\log T}{2\pi}\int_{-\beta}^\beta r(\alpha)F(\alpha)\dd\alpha=\sum_{0<\gamma,\gamma'\le T}w(\gamma-\gamma')\,\beta\Bigl(\frac{\sin(\beta v/2)}{\beta v/2}\Bigr)^2,\qquad v=(\gamma-\gamma')\log T.
\]
All terms on the right are non-negative, and the pairs with $\gamma=\gamma'$ contribute $\beta$ each; there are $\sum_{\text{distinct }\rho}m_\rho^2=\sum_{0<\gamma\le T}m_\rho$ of them. On the left, Theorem~\ref{thm:montgomerypc} (with $F$ even) gives
\[
\int_{-\beta}^\beta r(\alpha)F(\alpha)\dd\alpha=2\int_0^\beta\Bigl(1-\frac\alpha\beta\Bigr)\bigl((1+o(1))T^{-2\alpha}\log T+\alpha\bigr)\dd\alpha+o(1)=1+\frac{\beta^2}3+o(1),
\]
since $2\int_0^\beta(1-\frac\alpha\beta)T^{-2\alpha}\log T\dd\alpha\to1$. With $N(T)\sim\frac T{2\pi}\log T$, $\sum m_\rho\le(\frac1\beta+\frac\beta3+o(1))N(T)$ for every $\beta<1$, and letting $\beta\to1$ gives $\frac43$. If $N_s$ zeros are simple, the others have $m_\rho\ge2$, so $\sum m_\rho\ge N_s+2(N(T)-N_s)$, and $N_s\ge2N(T)-\sum m_\rho\ge(\frac23-o(1))N(T)$.
\end{proof}

It is conjectured that all zeros are simple, and the proportion $\frac23$ has been pushed a little further under the Riemann hypothesis, to about $0.70$, by other methods. Unconditionally, Theorem~\ref{thm:proportion} gives more than $40\%$ simple zeros on the line.

This book began with zeros that stray from the line, and ends with zeros on it. In between, primes were counted in short intervals and in progressions, sieved, summed in threes, packed into bounded gaps, arranged in progressions of every length, and compared with random models at every turn. The models always win in the end, but each victory needed a new idea, and the deepest of them, the Riemann hypothesis and Montgomery's conjecture, are still waiting for theirs.

\begin{think}
Theorem~\ref{thm:montgomerypc} stops at $\alpha<1$ because the explicit formula involves primes up to $x=T^\alpha$, and the mean value of $\sum_{n\le x}a_nn^{-it}$ over $[0,T]$ is governed by its diagonal only when $x$ is smaller than $T$. Where exactly in the proof would $x>T$ break the argument, and what would one need to know about pairs of primes $p,p'$ with $|p-p'|$ small to go further?
\end{think}

\section{Exercises}

\begin{exercise}[$\star$]\label{ex:zreal}
Let $\vartheta(t)=\arg\bigl(\pi^{-it/2}\Gamma(\frac14+\frac{it}2)\bigr)$, defined continuously with $\vartheta(0)=0$, and $Z(t)=e^{i\vartheta(t)}\zeta(\frac12+it)$. Show that $\Xi(t)=-\frac12(t^2+\frac14)\pi^{-1/4}|\Gamma(\frac14+\frac{it}2)|\,Z(t)$. Deduce that $Z$ is real, that $|Z(t)|=|\zeta(\frac12+it)|$, and that the zeros of odd order of $\zeta$ on the critical line are exactly the points where $Z$ changes sign.
\end{exercise}

\begin{exercise}[$\star$]\label{ex:sumgamma}
Assume the Riemann hypothesis. (a) Take $x=1$ in Lemma~\ref{lem:montexplicit} to show that $\sum_\gamma\frac1{1+(t-\gamma)^2}=\frac12\log\tau+O(1)$. (b) Deduce that at most $\log\tau+O(1)$ ordinates lie in $[t-1,t+1]$, and compare with the average number $\frac1\pi\log\frac t{2\pi}$.
\end{exercise}

\begin{exercise}[$\star$]\label{ex:fejervalue}
For $\beta>0$ and real $v\ne0$, show that $\int_{-\beta}^\beta(1-\frac{|\alpha|}\beta)e^{i\alpha v}\dd\alpha=\beta\bigl(\frac{\sin(\beta v/2)}{\beta v/2}\bigr)^2$, and that $\int_{-\beta}^\beta(1-\frac{|\alpha|}\beta)|\alpha|\dd\alpha=\frac{\beta^2}3$.
\end{exercise}

\begin{exercise}[$\star\star$]\label{ex:meansqsigma}
Let $\frac12<\sigma\le1$. Show that $\frac1T\int_0^T|\zeta(\sigma+it)|^2\dd t\to\zeta(2\sigma)$ as $T\to\infty$.
\end{exercise}

\begin{exercise}[$\star\star$]\label{ex:distinctzeros}
Assume the Riemann hypothesis. Show that the number of \emph{distinct} zeros with $0<\gamma\le T$ is at least $(\frac34-o(1))N(T)$.
\end{exercise}

\begin{exercise}[$\star\star\star$]\label{ex:hilbertopt}
Show that the constant $\pi$ in Theorem~\ref{thm:hilbert} is optimal: for $0<\beta<\frac12$ and $x_n=e(n\beta)$, show that $\frac1N\sum_{m\ne n\le N}\frac{x_m\overline{x_n}}{m-n}\to i\pi(1-2\beta)$ as $N\to\infty$.
\end{exercise}

\section{Hints}

\begin{hint}{ex:zreal}
Write $s(s-1)=-(t^2+\frac14)$ and $\pi^{-s/2}\Gamma(\frac s2)=\pi^{-1/4}\,\pi^{-it/2}\Gamma(\frac14+\frac{it}2)$. For the sign changes, compare orders of vanishing of $\Xi$ and of $\zeta(\frac12+it)$.
\end{hint}

\begin{hint}{ex:sumgamma}
At $x=1$ the first sum of primes is empty, and the second is $\sum_n\Lambda(n)n^{-3/2-it}$, bounded by $-\frac{\zeta'}\zeta(\frac32)$. For (b), each ordinate in $[t-1,t+1]$ contributes at least $\frac12$.
\end{hint}

\begin{hint}{ex:fejervalue}
The integrand is even in $\alpha$ up to conjugation; compute $2\int_0^\beta(1-\frac\alpha\beta)\cos(\alpha v)\dd\alpha$ by one integration by parts, and use $1-\cos y=2\sin^2\frac y2$.
\end{hint}

\begin{hint}{ex:meansqsigma}
Repeat the proofs of Proposition~\ref{prop:meansqpoly} and Theorem~\ref{thm:hardylittlewood} with $(mn)^{-\sigma}$ in place of $(mn)^{-1/2}$ and $x_n=n^{1/2-\sigma+iT}$ in Hilbert's inequality. Every off-diagonal bound becomes $O(N^{2-2\sigma})=o(T)$.
\end{hint}

\begin{hint}{ex:distinctzeros}
If $D$ distinct zeros carry the multiplicities $m_1,\dots,m_D$, then $(\sum m_j)^2\le D\sum m_j^2$.
\end{hint}

\begin{hint}{ex:hilbertopt}
Group the terms by $k=m-n$: the sum is $2i\sum_{k=1}^{N-1}(N-k)\frac{\sin2\pi k\beta}k$. Use $\sum_{k\ge1}\frac{\sin ky}k=\frac{\pi-y}2$ for $0<y<2\pi$ and Cesàro summation.
\end{hint}

\section{Solutions}

\begin{solution}{ex:zreal}
\idea All the factors in $\xi(\frac12+it)$ except $\zeta$ are explicit, and only the argument of the Gamma factor is not real: that argument is exactly what $\vartheta$ removes.
\writeup With $s=\frac12+it$, $\frac12s(s-1)=-\frac12(t^2+\frac14)$ and $\pi^{-s/2}\Gamma(\frac s2)=\pi^{-1/4}|\Gamma(\frac14+\frac{it}2)|e^{i\vartheta(t)}$, since $|\pi^{-it/2}|=1$. So $\Xi(t)=-\frac12(t^2+\frac14)\pi^{-1/4}|\Gamma(\frac14+\frac{it}2)|e^{i\vartheta(t)}\zeta(s)$, which is the formula. The factor $c(t)=-\frac12(t^2+\frac14)\pi^{-1/4}|\Gamma(\frac14+\frac{it}2)|$ is real, smooth and negative, and $\Xi$ is real, so $Z=\Xi/c$ is real, and $|Z|=|\zeta(s)|$. A zero of order $m$ of $t\mapsto\zeta(\frac12+it)$ at $t_0$ is a zero of order $m$ of $\Xi$ and of $Z$, which are real analytic; such a function changes sign at $t_0$ if and only if $m$ is odd.
\end{solution}

\begin{solution}{ex:sumgamma}
\idea With $x=1$ the oscillating factor $x^{i\gamma}$ disappears and the lemma counts zeros near $t$.
\writeup (a) For $x=1$ the sum over $n\le x$ reduces to $n=1$, where $\Lambda(1)=0$, and $|\sum_{n>1}\Lambda(n)n^{-3/2-it}|\le\sum_n\Lambda(n)n^{-3/2}=-\frac{\zeta'}\zeta(\frac32)$. The remaining terms are $\log\tau+O(1)$, so $2\sum_\gamma\frac1{1+(t-\gamma)^2}=\log\tau+O(1)$.
(b) Each ordinate with $|\gamma-t|\le1$ contributes at least $\frac12$ to the sum in (a), and the others contribute positive amounts. So the number of such ordinates is at most $2(\frac12\log\tau+O(1))=\log\tau+O(1)$. The average number in an interval of length $2$ at height $t$ is $2\cdot\frac1{2\pi}\log\frac t{2\pi}$ (Volume~6, Theorem~6.9), so the bound is $\pi$ times the average: no interval of length $2$ contains more than about $\pi$ times its share.
\end{solution}

\begin{solution}{ex:fejervalue}
\idea A triangle is the convolution of two rectangles, so its transform is the square of the transform of a rectangle.
\writeup The imaginary part of the integrand is odd in $\alpha$, so the integral is $2\int_0^\beta(1-\frac\alpha\beta)\cos(\alpha v)\dd\alpha$. Integrating by parts, $\int_0^\beta(1-\frac\alpha\beta)\cos(\alpha v)\dd\alpha=\bigl[(1-\frac\alpha\beta)\frac{\sin\alpha v}v\bigr]_0^\beta+\frac1\beta\int_0^\beta\frac{\sin\alpha v}v\dd\alpha=\frac{1-\cos\beta v}{\beta v^2}$. So the integral is $\frac{2(1-\cos\beta v)}{\beta v^2}=\frac{4\sin^2(\beta v/2)}{\beta v^2}=\beta\bigl(\frac{\sin(\beta v/2)}{\beta v/2}\bigr)^2$. Finally $2\int_0^\beta(\alpha-\frac{\alpha^2}\beta)\dd\alpha=2(\frac{\beta^2}2-\frac{\beta^2}3)=\frac{\beta^2}3$.
\end{solution}

\begin{solution}{ex:meansqsigma}
\idea To the right of the critical line the diagonal $\sum n^{-2\sigma}$ converges, and the off-diagonal is a power of $N$ smaller than $T$, so the same proof gives an asymptotic with a constant instead of a logarithm.
\writeup Let $x=T+3$, $N=\lfloor x\rfloor$ and $S_\sigma(t)=\sum_{n\le N}n^{-\sigma-it}$. As in Proposition~\ref{prop:meansqpoly}, $\int_0^T|S_\sigma|^2=T\sum_{n\le N}n^{-2\sigma}+R$, where $R$ is $\frac1i\sum_{m\ne n}(mn)^{-\sigma}\frac{(n/m)^{iT}}{\log(n/m)}$, the part with $-1$ cancelling by symmetry. Split $\frac1{\log(n/m)}=\frac{\sqrt{mn}}{n-m}+k(\log\frac nm)$. The first piece is $\frac1i\sum_{m\ne n}\frac{x_n\overline{x_m}}{n-m}$ with $x_n=n^{1/2-\sigma+iT}$, at most $\pi\sum_{n\le N}n^{1-2\sigma}=O(N^{2-2\sigma})$ by Hilbert's inequality. In the second, the terms with $m<n\le2m$ are at most $\frac{n-m}{20m^{1+2\sigma}}$, adding up to $O(\sum_mm^{1-2\sigma})=O(N^{2-2\sigma})$, and those with $n>2m$ add up to at most $\frac2{\log2}(\sum_{m\le N}m^{-\sigma})^2=O(N^{2-2\sigma})$. Since $\sigma>\frac12$, $R=o(T)$.
By Lemma~\ref{lem:zetalengtht}, $\zeta(\sigma+it)=S_\sigma(t)+r(t)$ with $|r(t)|\le\frac{x^{1-\sigma}}{\max(\sigma,t)}+4x^{-\sigma}$, so $\int_0^T|r|^2=O(x^{2-2\sigma}+Tx^{-2\sigma})=o(T)$. By Cauchy--Schwarz, $\int_0^T|\zeta(\sigma+it)|^2=T\sum_{n\le N}n^{-2\sigma}+o(T)$, and $\sum_{n\le N}n^{-2\sigma}\to\zeta(2\sigma)$.
\end{solution}

\begin{solution}{ex:distinctzeros}
\idea Corollary~\ref{cor:simplezeros} bounds the second moment of the multiplicities; few distinct zeros would force a large second moment.
\writeup Let the distinct zeros with $0<\gamma\le T$ have multiplicities $m_1,\dots,m_D$. Then $\sum m_j=N(T)$, and $\sum_{0<\gamma\le T}m_\rho=\sum_jm_j^2\le(\frac43+o(1))N(T)$ by Corollary~\ref{cor:simplezeros}. By Cauchy--Schwarz, $N(T)^2=(\sum m_j)^2\le D\sum m_j^2\le D(\frac43+o(1))N(T)$, so $D\ge(\frac34-o(1))N(T)$.
\end{solution}

\begin{solution}{ex:hilbertopt}
\idea For the extremal vectors the quadratic form is close to a multiple of $\sum_{k\ne0}\frac{e(k\beta)}k$, which is the Fourier series of a sawtooth, and the sawtooth jumps by $\pi$ at $0$.
\writeup Grouping by $k=m-n$, each $k$ with $0<|k|<N$ occurs $N-|k|$ times, so the sum is $\sum_{0<|k|<N}(N-|k|)\frac{e(k\beta)}k=2i\sum_{k=1}^{N-1}(N-k)\frac{\sin2\pi k\beta}k$. Divided by $N$, this is $2i$ times the Cesàro mean of the partial sums of $\sum_{k\ge1}\frac{\sin2\pi k\beta}k$. That series converges to $\frac{\pi-2\pi\beta}2$ (the Fourier series of the sawtooth, Volume~5, Chapter~8), and Cesàro means of a convergent series have the same limit. So the normalised sum tends to $2i\cdot\frac{\pi(1-2\beta)}2=i\pi(1-2\beta)$, while $\frac1N\sum|x_n|^2=1$. Letting $\beta\to0$, the ratio approaches $\pi$, and no smaller constant works.
\end{solution}

\section*{Sources and further reading}

\begin{itemize}
\item \emph{Hardy's theorem.} G.~H.~Hardy, ``Sur les zéros de la fonction $\zeta(s)$ de Riemann'', \emph{C.\ R.\ Acad.\ Sci.\ Paris} 158 (1914), 1012--1014. Our proof follows E.~C.~Titchmarsh, \emph{The Theory of the Riemann Zeta-Function}, 2nd ed., revised by D.~R.~Heath-Brown, Oxford, 1986, Section~10.2, and its formula (2.16.1). The postcard: G.~Pólya, ``Some mathematicians I have known'', \emph{Amer.\ Math.\ Monthly} 76 (1969), 746--753.
\item \emph{Proportions.} G.~H.~Hardy and J.~E.~Littlewood, ``The zeros of Riemann's zeta-function on the critical line'', \emph{Math.\ Z.}\ 10 (1921), 283--317. A.~Selberg, ``On the zeros of Riemann's zeta-function'', \emph{Skr.\ Norske Vid.\ Akad.\ Oslo} I (1942), no.~10. N.~Levinson, ``More than one third of zeros of Riemann's zeta-function are on $\sigma=1/2$'', \emph{Adv.\ Math.}\ 13 (1974), 383--436. J.~B.~Conrey, ``More than two fifths of the zeros of the Riemann zeta function are on the critical line'', \emph{J.\ reine angew.\ Math.}\ 399 (1989), 1--26. K.~Pratt, N.~Robles, A.~Zaharescu and D.~Zeindler, ``More than five-twelfths of the zeros of $\zeta$ are on the critical line'', \emph{Res.\ Math.\ Sci.}\ 7 (2020), arXiv:1802.10521. D.~Platt and T.~Trudgian, ``The Riemann hypothesis is true up to $3\cdot10^{12}$'', \emph{Bull.\ London Math.\ Soc.}\ 53 (2021), 792--797.
\item \emph{Mean values.} G.~H.~Hardy and J.~E.~Littlewood, ``Contributions to the theory of the Riemann zeta-function and the theory of the distribution of primes'', \emph{Acta Math.}\ 41 (1918), 119--196. A.~E.~Ingham, ``Mean-value theorems in the theory of the Riemann zeta-function'', \emph{Proc.\ London Math.\ Soc.}\ (2) 27 (1926), 273--300. I.~Schur, ``Bemerkungen zur Theorie der beschränkten Bilinearformen mit unendlich vielen Veränderlichen'', \emph{J.\ reine angew.\ Math.}\ 140 (1911), 1--28. H.~L.~Montgomery and R.~C.~Vaughan, ``Hilbert's inequality'', \emph{J.\ London Math.\ Soc.}\ (2) 8 (1974), 73--82. The approximate functional equation is Titchmarsh's Theorem~4.13.
\item \emph{Computing.} C.~L.~Siegel, ``Über Riemanns Nachlaß zur analytischen Zahlentheorie'', \emph{Quellen und Studien zur Geschichte der Math.\ Astr.\ Physik} B~2 (1932), 45--80. H.~M.~Edwards, \emph{Riemann's Zeta Function}, Academic Press, 1974, Chapter~7. A.~M.~Odlyzko and A.~Schönhage, ``Fast algorithms for multiple evaluations of the Riemann zeta function'', \emph{Trans.\ Amer.\ Math.\ Soc.}\ 309 (1988), 797--809.
\item \emph{Pair correlation.} H.~L.~Montgomery, ``The pair correlation of zeros of the zeta function'', in \emph{Analytic Number Theory}, Proc.\ Sympos.\ Pure Math.\ 24, Amer.\ Math.\ Soc., 1973, 181--193. D.~A.~Goldston and H.~L.~Montgomery, ``Pair correlation of zeros and primes in short intervals'', in \emph{Analytic Number Theory and Diophantine Problems}, Progr.\ Math.\ 70, Birkhäuser, 1987, 183--203. D.~A.~Goldston, ``Notes on pair correlation of zeros and prime numbers'', arXiv:math/0412313. The tea at Princeton: K.~Sabbagh, \emph{Dr Riemann's Zeros}, Atlantic Books, 2002, and Volume~8, Chapter~13.
\end{itemize}
